\chapter{Минимальные гамильтоновы данные центрального веса синглета и триплета}
\label{ch:v8-baryon-c0-singlet-triplet-central-weight-minimal-hamiltonian-data-gate}

Литературный проход уточняет постановку: максимум энтропии требует заранее
заданных ограничений, Gibbs-вариация требует гамильтониана, а detailed
balance определяется относительно уже выбранного верного состояния.
Поэтому минимизируем не число параметров записи, а число физически
различимых данных.

\section{Факторизация допустимого гамильтониана}

По лемме Шура всякий $SO(3)$-инвариантный гамильтониан на
$\mathbf1\oplus\mathbf3$ имеет вид
\begin{equation}
 h=\varepsilon_1P_{\mathbf1}+\varepsilon_3P_{\mathbf3}
 =\varepsilon_1I_4+\Delta P_{\mathbf3},
 \qquad \Delta=\varepsilon_3-\varepsilon_1.
 \label{eq:v8-baryon-c0-minimal-central-hamiltonian-form}
\end{equation}
Добавление общего нуля энергии не меняет состояния:
\begin{equation}
 h\sim h+bI_4,
 \qquad
 \frac{e^{-\beta(h+bI_4)}}{\Tr e^{-\beta(h+bI_4)}}
 =\frac{e^{-\beta h}}{\Tr e^{-\beta h}}.
 \label{eq:v8-baryon-c0-central-hamiltonian-shift-quotient}
\end{equation}
После этой факторизации остаётся одна энергетическая разность. Для
равновесного веса существенна только безразмерная комбинация
\begin{equation}
 \theta=\beta\Delta.
 \label{eq:v8-baryon-c0-central-dimensionless-gap}
\end{equation}

Её Gibbs-состояние равно
\begin{equation}
 \rho_\theta=p(\theta)P_{\mathbf1}
 +\frac{1-p(\theta)}{3}P_{\mathbf3},
 \qquad
 p(\theta)=\frac{1}{1+3e^{-\theta}},
 \qquad r(\theta)=e^{-\theta}.
 \label{eq:v8-baryon-c0-central-gap-gibbs-state}
\end{equation}
Отображение является биекцией
\begin{equation}
 \theta\in\mathbb R\longleftrightarrow p\in(0,1),
 \qquad
 \theta(p)=\log\frac{3p}{1-p}.
 \label{eq:v8-baryon-c0-central-gap-weight-bijection}
\end{equation}
Его монотонность невырождена во внутренности:
\begin{equation}
 \frac{dp}{d\theta}=p(1-p)>0.
 \label{eq:v8-baryon-c0-central-gap-weight-monotonicity}
\end{equation}
Тем самым один вещественный параметр необходим и достаточен для выбора
центрального веса. Два точных ориентира имеют вид
\begin{equation}
 \theta=0\Longleftrightarrow(p,r)=\left(\frac14,1\right),
 \qquad
 \theta=\log3\Longleftrightarrow(p,r)=\left(\frac12,\frac13\right).
 \label{eq:v8-baryon-c0-central-gap-two-weight-witnesses}
\end{equation}

\section{Gibbs-вариация как проверка достаточности}

На центральном симплексе введём безразмерную свободную энергию
\begin{equation}
 \Phi_\theta(p)=(1-p)\theta-S(p),
 \qquad
 S(p)=-p\log p-(1-p)\log\frac{1-p}{3}.
 \label{eq:v8-baryon-c0-central-gap-free-energy-functional}
\end{equation}
Её производные равны
\begin{equation}
 \Phi_\theta'(p)=\log\frac{3p}{1-p}-\theta,
 \qquad
 \Phi_\theta''(p)=\frac{1}{p(1-p)}>0.
 \label{eq:v8-baryon-c0-central-gap-free-energy-derivatives}
\end{equation}
Поэтому $p(\theta)$ является единственным минимумом. Более сильная точная
форма есть Gibbs-вариационное тождество
\begin{equation}
 \Phi_\theta(p)-\Phi_\theta(p(\theta))
 =D\!\left(\rho_p\middle\|\rho_\theta\right)\ge0,
 \label{eq:v8-baryon-c0-central-gap-relative-entropy-identity}
\end{equation}
причём равновесное значение равно
\begin{equation}
 \Phi_\theta(p(\theta))=-\log\left(1+3e^{-\theta}\right).
 \label{eq:v8-baryon-c0-central-gap-equilibrium-free-energy}
\end{equation}
Это доказывает достаточность $\theta$, но не выводит его значение.

\section{Что равновесие не идентифицирует}

Преобразование
\begin{equation}
 (\beta,\Delta)\longmapsto
 \left(\frac{\beta}{a},a\Delta\right),
 \qquad a>0,
 \label{eq:v8-baryon-c0-central-gap-temperature-energy-orbit}
\end{equation}
сохраняет $\theta$, $p$, $r$ и всю равновесную матрицу плотности. Поэтому
равновесный вес не различает температуру и физическую энергетическую щель.
Кроме того, Davies-скорости требуют спектральной плотности среды и не
следуют из одной $\rho_\theta$.

Минимальный контракт распадается на три слоя:
\begin{equation}
 \underbrace{\theta=\beta\Delta}_{\text{центральный вес}},
 \qquad
 \underbrace{\beta\ \text{или}\ \Delta_{\rm phys}}_{\text{абсолютная энергия}},
 \qquad
 \underbrace{\kappa_{\rm bath}>0}_{\text{время релаксации}}.
 \label{eq:v8-baryon-c0-central-gap-three-data-layers}
\end{equation}
Для текущего вопроса нужен только первый слой. Он одномерен, но всё ещё не
получен из родителя:
\begin{equation}
 \dim\mathcal S_{\rm weight}^{\min}=1,
 \qquad
 N_{\rm derived}(\theta)=0/1.
 \label{eq:v8-baryon-c0-central-gap-minimal-data-ledger}
\end{equation}

\begin{theorem}[Минимальность безразмерной центральной щели]
С точностью до общего сдвига энергии всякий $SO(3)$-инвариантный
гамильтониан типа $\mathbf1\oplus\mathbf3$ задаётся одной щелью $\Delta$.
При фиксированной температуре, а в безразмерной формулировке без неё,
единственное необходимое и достаточное данное для выбора $p\in(0,1)$ есть
$\theta=\beta\Delta$. Gibbs-вариационный функционал имеет при этом
единственный минимум $p(\theta)$.
\end{theorem}

\begin{nogocriterion}[Равновесный вес не задаёт энергию и скорость]
Нельзя извлечь из одного $p$ раздельные значения $\beta$ и
$\Delta_{\rm phys}$: они лежат на масштабной орбите с постоянным
$\beta\Delta$. Нельзя также считать $r=e^{-\theta}$ абсолютной скоростью
релаксации без корреляционной нормы среды.
\end{nogocriterion}

\begin{protocol}\begin{enumerate}
\item допустимый центральный гамильтониан: \textbf{$\varepsilon I_4+\Delta P_{\mathbf3}$};
\item общий сдвиг энергии физически различим: \textbf{нет};
\item минимальный селектор веса: \textbf{$\theta=\beta\Delta$};
\item пространство селектора: \textbf{$\mathbb R$};
\item отображение $\theta\mapsto p$: \textbf{биективно};
\item Gibbs-минимум: \textbf{единственный};
\item $\beta$ и $\Delta$ отдельно идентифицируются: \textbf{нет};
\item скорость релаксации определяется равновесием: \textbf{нет};
\item размерность минимальных данных веса: \textbf{$1$};
\item выведено текущим родителем: \textbf{$0/1$};
\item следующий гейт: \textbf{происхождение центральной щели из родительского действия}.
\end{enumerate}\end{protocol}

Литературная опора: E.~T.~Jaynes, \emph{Information Theory and Statistical
Mechanics}; E.~B.~Davies, \emph{Markovian Master Equations};
F.~Fagnola и V.~Umanit\`a, \emph{Generators of Detailed Balance Quantum
Markov Semigroups}; Gibbs-вариационная формула в arXiv:2111.10282.

Машинный сертификат сохранён в
\path{s2t_v8_baryon_c0_singlet_triplet_central_weight_minimal_hamiltonian_data_gate_results.json}.